\documentclass[sigconf,nonacm]{acmart}

% Pacotes complementares
\usepackage[utf8]{inputenc}
\usepackage{booktabs}
\usepackage{graphicx}

% Metadados do artigo
\title{Título do Trabalho Científico: Uma Proposta Aplicada com Impacto Comunitário}

% Autores e Afiliação Institucional
\author{Estudante 1}
\affiliation{%
  \institution{Instituto Federal de Santa Catarina (IFSC)}
  \city{Garopaba}
  \state{SC}
  \country{Brasil}
}
\email{estudante1@aluno.ifsc.edu.br}

\author{Estudante 2}
\affiliation{%
  \institution{Instituto Federal de Santa Catarina (IFSC)}
  \city{Garopaba}
  \state{SC}
  \country{Brasil}
}
\email{estudante2@aluno.ifsc.edu.br}

\author{André Luiz Silva de Moraes}
\affiliation{%
  \institution{Instituto Federal de Santa Catarina (IFSC)}
  \city{Garopaba}
  \state{SC}
  \country{Brasil}
}
\email{andre.moraes@ifsc.edu.br}

\author{Adalberto Teodosio Tabalipa}
\affiliation{%
  \institution{Instituto Federal de Santa Catarina (IFSC)}
  \city{Garopaba}
  \state{SC}
  \country{Brasil}
}
\email{adalberto.tabalipa@ifsc.edu.br}

\begin{document}

\begin{abstract}
Este trabalho apresenta um estudo experimental e extensionista desenvolvido no âmbito da Unidade Curricular de Iniciação Científica e Extensionista do CST em Sistemas para a Internet do IFSC Câmpus Garopaba. O artigo aborda a investigação aplicada de tecnologias computacionais e seus impactos na comunidade regional.
\end{abstract}

\keywords{Iniciação Científica, Extensão Universitária, Sistemas para Internet, COTB}

\maketitle

\section{Introdução}\label{sec:introducao}
O avanço das tecnologias de informação e comunicação tem proporcionado novas oportunidades de extensão universitária e desenvolvimento regional~\cite{cotb_exemplo2026}. Neste contexto, o presente trabalho investiga soluções aplicadas para desafios tecnológicos contemporâneos.

Como demonstrado em estudos prévios na literatura especializada~\cite{kalman1960, haykin2014}, o rigor metodológico na coleta de dados é essencial para a reprodutibilidade dos experimentos.

\section{Metodologia e Resultados}\label{sec:metodologia}
A abordagem metodológica divide-se em três etapas principais: (i) revisão sistemática da literatura e mapeamento do estado da arte; (ii) prototipação experimental com scripts Python de análise estatística; e (iii) validação das métricas junto à comunidade externa.

A Figura~\ref{fig:desempenho} ilustra a arquitetura conceitual e o fluxo de dados implementado durante a pesquisa.

\begin{figure}[htbp]
  \centering
  \includegraphics[width=\columnwidth]{figura_exemplo.png}
  \caption{Visão geral do pipeline de dados e arquitetura da solução proposta.}
  \label{fig:desempenho}
\end{figure}

Os dados quantitativos obtidos durante a experimentação encontram-se sumarizados na Tabela~\ref{tab:resultados}, formatada conforme as diretrizes internacionais da ACM.

\begin{table}[htbp]
  \caption{Métricas de desempenho e tempo médio de resposta}
  \label{tab:resultados}
  \centering
  \begin{tabular}{lcr}
    \toprule
    \textbf{Ferramenta Avaliada} & \textbf{Amostras (N)} & \textbf{Taxa de Sucesso (\%)} \\
    \midrule
    Script Python Automatizado   & 50                    & 98,4 \\
    Auditoria Manual Heurística  & 50                    & 86,2 \\
    Coleta Baseada em API Web    & 50                    & 94,8 \\
    \bottomrule
  \end{tabular}
\end{table}

\section{Conclusão e Trabalhos Futuros}\label{sec:conclusao}
Os resultados preliminares demonstram a viabilidade técnica da abordagem. Como trabalhos futuros, planeja-se a ampliação do conjunto de dados e a realização de novos testes de campo na região da AMUREL e Grande Florianópolis.

\section*{Agradecimentos}
Os autores agradecem ao Instituto Federal de Santa Catarina (IFSC) pelo apoio institucional e aos docentes e colegas do curso pelo suporte acadêmico.

\section*{Declaração de Uso de Inteligência Artificial}
Os autores declaram que ferramentas de Inteligência Artificial generativa foram utilizadas estritamente para auxílio na revisão ortográfica e formatação de trechos em \LaTeX, mantendo total responsabilidade humana pela integridade científica do conteúdo.

% Referências bibliográficas via BibTeX
\bibliographystyle{ACM-Reference-Format}
\bibliography{referencias}

\end{document}
